\chapter{Introducción}\label{cap:introduccion}

El propósito de este trabajo es resolver la complejidad inherente a la logística de distribución mediante el uso de bases de datos de grafos. En lugar de utilizar modelos relacionales tradicionales que sufren con consultas recursivas pesadas al intentar modelar redes, he optado por ArangoDB para modelar almacenes como nodos y rutas como aristas, permitiendo una adaptabilidad inmediata ante cambios en el entorno.

La metodología seguida ha sido iterativa: partiendo de un diseño de esquema básico, progresé hasta lograr una arquitectura modular robusta, enfrentándome a retos técnicos tanto en la lógica de resolución de rutas óptimas mediante AQL como en la orquestación del despliegue en contenedores Docker, los cuales se detallan en los posteriores capítulos de implementación.